\documentclass[12pt]{article}
\usepackage[margin=1in]{geometry} 
\usepackage{amsmath,amsthm,amssymb}
\usepackage[T2A]{fontenc}
\usepackage[russian]{babel}
\usepackage{lipsum}
\usepackage{graphicx}
\usepackage{listings}
\usepackage{xcolor}
\usepackage{float}

\newcommand{\N}{\mathbb{N}}
\newcommand{\Z}{\mathbb{Z}}
\newcommand{\forceindent}{\leavevmode{\parindent=2em\indent}}
% \renewcommand{\baselinestretch}{1.5} 
\setlength{\parindent}{0pt}

\title{Численные методы решения стационарного уравнения теплопроводности. Отчет}
\author{Новикова Вера, 301 группа}
\date{Вариант 4.3}

\begin{document}
\maketitle

\section*{Постановка задачи}
В данной работе реализовано численное решение краевой задачи для одномерного стационарного уравнения диффузии:
\[
\begin{cases}
    -\frac{d}{dx} \bigg(p(x) \frac{du}{dx}\bigg) +q(x)u(x)= f(x), & x \in [0, L], \text{ } L = \pi\\
    -p(0)\frac{du}{dx}(0) + \kappa_0 u(0) = g_0, & \kappa_0 = 1, g_0 = 1 \\ 
    u(L) = g_1, & g_1 = 1 
\end{cases}
\]

\begin{center}
\begin{minipage}{0.45\textwidth}
\begin{equation*}
p(x) =
\begin{cases} 
1 + \frac{1}{4}\sin(x), & x \in [0; \frac{\pi}{3}), \\
1 + \frac{1}{3}\sin(x), & x \in \left(\frac{\pi}{3}, \frac{2\pi}{3}\right), \\
1 + \frac{1}{2}\sin(x), & x \in \left(\frac{2\pi}{3}, \pi\right]
\end{cases}
\end{equation*}
\end{minipage}
\hfill
\begin{minipage}{0.45\textwidth}
\begin{equation*}
q(x) =
\begin{cases} 
1, & x \in [0; \frac{\pi}{3}), \\
2, & x \in \left(\frac{\pi}{3}, \frac{2\pi}{3}\right), \\
1, & x \in \left(\frac{2\pi}{3}, \pi\right]
\end{cases}
\end{equation*}
\end{minipage}
\end{center}
\[
f(x) = 1-cos(2x)
\]
Была построена консервативная схема на неравномерной расчетной сетке и приведен вывод ее погрешности аппроксимации, а также продемонстрирована сходимость метода с помощью метода Рунге.

\section*{Вывод порядка погрешности аппроксимации схемы на неравномерной сетке}
Решаем уравнение:
\[
-\frac{d}{dx} \bigg(p(x) \frac{du}{dx}\bigg) +q(x)u(x)= f(x)
\]
Перепишем его в виде:
\begin{equation}
-p(x)u''(x)-p'(x)u'(x)+q(x)u(x)-f(x)=0
\end{equation}
Пусть $\{x_i\}$ - узлы сетки. Обозначим 
\[
x_{i+\frac{1}{2}} = \frac{x_i+x_{i+1}}{2}, \text{ } x_{i-\frac{1}{2}} = \frac{x_{i-1}+x_{i}}{2}, \text{ } h_i=x_{i+1}-x_i, \text{ } h_{i-1}=x_i-x_{i-1}
\]
Проинтегрируем $(1)$ от $x_{i-1/2}$ до $x_{i+1/2}$:
\[
-p(x)\frac{du}{dx}\bigg|_{x_{i-1/2}}^{x_{i+1/2}} + \int_{x_{i-1/2}}^{x_{i+1/2}}\big[q(x)u(x)-f(x)\big]dx =0
\]
Если функция $q(x)$ разрывна, то при интегрировании возникнут проблемы. Для этого приблизим ее некоторым средним значением: $\overline{q_i} = \frac{1}{h}\int_{x_{i-1/2}}^{x_{i+1/2}}q(x)dx$. Тогда $\int qu \approx \overline{q}\int u$.

Ищем погрешность аппроксимации решения:
\[
-p_{i+1/2}\frac{u_{i+1}-u_i}{h_i}+p_{i-1/2}\frac{u_i-u_{i-1}}{h_{i-1}} + \frac{h_i+h_{i-1}}{2}\big[\overline{q_i}u_i-f_i\big]=0
\]
\[
-\frac{2}{h_i+h_{i-1}}\bigg[p_{i+1/2}\frac{u_{i+1}-u_i}{h_i} - p_{i-1/2}\frac{u_i-u_{i-1}}{h_{i-1}}\bigg]+q_iu_i-f_i=0
\]
Для этого разложим $u_{i\pm 1}$ и $p_{i\pm 1/2}$ в ряд Тейлора:
\[
u_{i + 1} = u_i+h_i\frac{du_i}{dx} + \frac{h_i^2}{2} \frac{d^2u_i}{dx^2} +  \frac{h_i^3}{6} \frac{d^3u_i}{dx^3} + \mathrm{O}(h_i^4)
\]
\[
u_{i - 1} = u_i-h_i\frac{du_i}{dx} + \frac{h_i^2}{2} \frac{d^2u_i}{dx^2} - \frac{h_i^3}{6} \frac{d^3u_i}{dx^3} + \mathrm{O}(h_i^4)
\]
\[
p_{i + 1/2} = p_i+\frac{h_i}{2}\frac{dp_i}{dx} + \frac{h_i^2}{8} \frac{d^2p_i}{dx^2} +  \frac{h_i^3}{48} \frac{d^3p_i}{dx^3} + \mathrm{O}(h_i^4)
\]
\[
p_{i - 1/2} = p_i-\frac{h_i}{2}\frac{dp_i}{dx} + \frac{h_i^2}{8} \frac{d^2p_i}{dx^2} - \frac{h_i^3}{48} \frac{d^3p_i}{dx^3} + \mathrm{O}(h_i^4)
\]
Обозначим $\overline{h_i} = \frac{h_i+h_{i-1}}{2}$, $h=\max\limits_{i}(h_i)$ 
Тогда оценим погрешность аппроксимации $\Psi$ как:
\[
\Psi = -\frac{1}{\overline{h_i}}\bigg[\big( p_i+\frac{h_i}{2}\frac{dp_i}{dx} + \frac{h_i^2}{8} \frac{d^2p_i}{dx^2} + \mathrm{O}(h^3)\big) \big(\frac{du_i}{dx}+\frac{h_i}{2}\frac{d^2u_i}{dx^2} +\frac{h_i^2}{6}\frac{d^3u_i}{dx^3} + \mathrm{O}(h^3)\big) - 
\]
\[
- \big( p_i-\frac{h_{i-1}}{2}\frac{dp_i}{dx} + \frac{h_{i-1}^2}{8} \frac{d^2p_i}{dx^2} + \mathrm{O}(h^3) \big) \big( \frac{du_i}{dx}-\frac{h_{i-1}}{2}\frac{d^2u_i}{dx^2} +\frac{h_{i-1}^2}{6}\frac{d^3u_i}{dx^3} + \mathrm{O}(h^3)\big)\bigg] + q_iu_i-f_i =
\]
\[
= -\frac{1}{\overline{h_i}}\bigg[ p_i\big(\overline{h_i}\frac{d^2u_i}{dx^2} + \frac{h_i^2-h_{i-1}^2}{6} \frac{d^3u_i}{dx^3} + \mathrm{O}(h^3)\big) + \frac{dp_i}{dx}\big( \overline{h_i}\frac{du_i}{dx} + \frac{h_i^2-h_{i-1}^2}{4} \frac{d^2u_i}{dx^2} + \mathrm{O}(h^3)\big) +
\]
\[
+ \frac{d^2p_i}{dx^2}\big( \frac{h_i^2-h_{i-1}^2}{8}\frac{du_i}{dx} + \mathrm{O}(h^3)\big)\bigg]+ q_iu_i - f_i = -p_i\frac{d^2u_i}{dx^2}-p_i\frac{h_i-h_{i-1}}{3}\frac{d^3u_i}{dx^3} -
\]
\[
- \frac{dp_i}{dx} \frac{du_i}{dx} -\frac{dp_i}{dx}\frac{h_i-h_{i-1}}{2} \frac{d^2u_i}{dx^2} - \frac{d^2p_i}{dx^2} \frac{h_i-h_{i-1}}{4} \frac{du_i}{dx} +q_iu_i -f_i+ \mathrm{O}(h^2) =
\]
\[
= \{(1)\} = (h_i-h_{i-1})(-\frac{p_i}{3}\frac{d^3u_i}{dx^3} - \frac{1}{2} \frac{dp_i}{dx} \frac{d^2u_i}{dx^2} - \frac{1}{4} \frac{dp_i}{dx} \frac{du_i}{dx}) + \mathrm{O}(h^2) = C(h_i-h_{i-1}) + \mathrm{O}(h^2)
\]
Если $h_i-h_{i-1} = \mathrm{O}(h^2)$, то общая погрешность аппроксимации порядка 2.

Ищем погрешность аппроксимации граничного условия:
\[
-p(0)\frac{du}{dx}(0) + \kappa_0 u(0) = g_0
\]
Проинтегрируем (1) от $x_0$ до $x_{1/2}$:
\[
-p(x)\frac{du}{dx}\bigg|_{x_0}^{x_{1/2}} + \int_{x_0}^{x_{1/2}}[q(x)u(x)-f(x)]dx
\]
С учетом:
\[
\int_0^{x_{1/2}} [q(x)u(x) - f(x)] \, dx \approx (q(0)u(0) - f(0)) \frac{h_1}{2}
\]
Получаем:
\[
-p_{1/2}u'_{1/2} + \frac{h_1}{2}\big(q(0)u(0) - f(0)\big) = -p(0)u'(0) = g_0-\kappa_0u(0)
\]
Разложим по Тейлору:
\[
-\big(p(0) + \frac{h_1}{2} p'(0) + \mathrm{O}(h^2)\big) \big(u'(0) + \frac{h_1}{2} u''(0) + \mathrm{O}(h^2) \big) + \frac{h_1}{2}\big(q(0)u(0) - f(0)\big)  = g_0^h-\kappa_0u(0)
\]
\[
-p(0)u'(0) - \frac{h_1}{2} p'(0) u'(0) - \frac{h_1}{2} p(0) u''(0) + \frac{h_1}{2} (q(0)u(0) - f(0))+ \kappa_0 u(0) + O(h^2) = g_0^h
\]

Погрешность
\[
\Psi = g_0 - g_0^h = g_0 + p(0)u'(0)  + \frac{h_1}{2} p'(0) u'(0)+ \frac{h_1}{2} p(0) u''(0) - \frac{h_1}{2} \big(q(0)u(0) - f(0)\big) - \kappa_0 u(0) + O(h^2)
\]

В силу граничных условий и равенства уравнения погрешность \(\Psi = O(h^2)\) (остальные слагаемые сокращаются).

Итак, для того, чтобы порядок аппроксимации схемы был равен 2, необходимо выполнение условия \( h_{i} - h_{i-1} = O(h^2) \), где $h=\max\limits_{i}(h_i)$.

\section*{Построение неравномерной сетки}
Обозначим $\phi(t): [0,1] \rightarrow [\alpha, \beta]$ - вспомогательная функция, по которой будем выбирать узлы сетки. На $\phi$ наложим следующие ограничения:
\begin{enumerate}
	\item $\phi(t) \in C^2[0, 1]$
	\item $\phi'(t) > 0$
\end{enumerate}
Покажем, что при $\phi(t) = \alpha + (\beta - \alpha) \frac{e^{ct}-1}{e^c-1}$ построенная неравномерная сетка $\{x_i\} = \{ \phi(t_i)\}, t_i=i\theta, \theta = \frac{1}{n}, i = \overline{1, n}$ удовлетворяет условию для 2 порядка аппроксимации схемы.
\[
\phi(t) \in C^2[0, 1], \quad \phi'(t) = (\beta - \alpha)\frac{ce^{ct}}{e^c-1} > 0
\]
\[
h_{i+1} = x_{i+1} - x_{i} = \phi(t_{i+1}) - \phi(t_i) = (\beta - \alpha) \frac{e^{c(i+1)\theta}-e^{ci\theta}}{e^c-1} = (\beta - \alpha) e^{ci\theta}\frac{e^{c\theta}-1}{e^c-1}
\]
\[
h_{i} = x_{i} - x_{i-1} = \phi(t_{i}) - \phi(t_{i-1}) = (\beta - \alpha) \frac{e^{ci\theta}-e^{c(i-1)\theta}}{e^c-1} = (\beta - \alpha) e^{ci\theta}\frac{1- e^{-c\theta}}{e^c-1}
\]
\[
\frac{h_{i+1} - h_i}{h_i^2} = \frac{\beta-\alpha}{e^c-1}e^{ci\theta}(e^{c\theta}-2+e^{-c\theta})\bigg(\frac{e^c-1}{(\beta-\alpha)e^{ci\theta}(1-e^{-c\theta})}\bigg)^2 = \frac{e^{-c\theta}(e^{c\theta}-1)^2(e^c-1)}{(\beta-\alpha)e^{ci\theta}(1-e^{-c\theta})^2}
\]
\[
\frac{h_{i+1} - h_i}{h_i^2} \rightarrow const, \text{ } \theta \rightarrow 0
\]
Итак, для построения сетки будем использовать вспомогательную функцию (положим $\alpha = 0, \beta = \pi, c = 1$):
\[
\phi(t)=\pi\frac{e^t-1}{e-1}
\]



Будем строить сетку так, чтобы все разрывы попадали в узлы.
%\begin{figure}[h]
    %\centering
    %\includegraphics[width=0.47\textwidth]{oblast.png} % Путь к файлу изображения
%\end{figure}

\section*{Код программы}
\lstset{
    language=C++, % Set language
    basicstyle=\ttfamily\small, % Monospace font
    keywordstyle=\color{blue}, % Keywords in blue
    commentstyle=\color{green}, % Comments in green
    stringstyle=\color{red}, % Strings in red
    numbers=left, % Line numbers on the left
    numberstyle=\tiny\color{gray}, % Line number style
    frame=single, % Border around the code
    breaklines=true, % Automatic line breaking
    tabsize=4 % Tab size
}
\begin{lstlisting}
int main()
{
    return 0;
}
\end{lstlisting}
\end{document}